\subsection{Homogeneous convex regularization}
\label{subsec:homogeneous-convex-regularization}
\index{Regularization!homogeneous fibre norms}
\index{Vector bundles!continuous convex norms}
\index{Plurisubharmonic functions!homogeneous regularization}
\index{Greene--Wu gluing}

Let $\pi:F\to X$ be a holomorphic vector bundle over a compact complex manifold. Write $F^\times$ for the complement of its zero section. A continuous real function $Q$ on $F$ is \emph{complex homogeneous of degree two} if
\[
 Q(\lambda v)=|\lambda|^2Q(v).
\]
We call $Q$ \emph{locally strongly convex in the fibres} if, in a local bundle frame and near every point of $F^\times$, there is a constant $c>0$ such that $Q(z,v)-c|v|^2$ is convex on sufficiently small real fibre balls. For continuous functions, strict plurisubharmonicity means that subtracting a sufficiently small positive multiple of the Euclidean squared norm in local total-space coordinates leaves a plurisubharmonic function. Both conditions are local and allow uniform positive margins on compact subsets of $F^\times$.

The smooth-norm obstruction also applies to continuous homogeneous functions with strict positivity margins. The local regularization is Wu's projective-coordinate convolution, written as an average of holomorphic linear shears. Compact gluing follows the argument of Greene and Wu. We give both steps to track strict plurisubharmonicity and strong fibre convexity; the exact attribution is recorded in \cref{sec:history-gauduchon}.

\begin{lemma}\label{lem:compact-sheaf-regularization}
Let $M$ be a compact smooth manifold and let $\mathcal S$ be a sheaf of germs of continuous real functions. Suppose that it is closed under maxima and has the following two properties.

First, if $f$ is a section near a compact set $K$, then $f+u$ is a section near $K$ whenever $u$ is smooth and sufficiently small in the $C^\infty$ topology on a neighbourhood of $K$. Second, every point has a neighbourhood $U$ on which sections admit semilocal smooth approximation: on a compact $K\Subset U$, a section can be approximated uniformly by smooth sections near $K$, with arbitrarily small errors in any prescribed finite collection of derivative norms on compact subsets where the section is already smooth.

Then every global section can be approximated uniformly on $M$ by smooth global sections.
\end{lemma}
\begin{proof}
We give the gluing step. Let $f$ be a global section, smooth near a compact set $L_1$, and let $L_2$ be a compact subset of one of the neighbourhoods $U$ in the statement. Choose relatively compact open sets $W_3,W_4$ with $L_1\subset W_3$, $f$ smooth near $\overline W_3$, and $L_2\subset W_4\Subset U$. Choose $0\leq\rho\leq1$, equal to one near $\overline W_3$, with compact support in the region where $f$ is smooth. Put $V=W_3\cup W_4$. A smooth function $-1\leq\psi\leq1$ can be chosen equal to one near $L_1\cup L_2$ and equal to minus one near $M\setminus V$.

By stability, for sufficiently small $\xi>0$ the function $F=f+\xi\psi$ is still a global section. Approximate $F$ on $\overline W_4$ by a smooth section $g$, with uniform error less than $\delta<\xi/2$. Require also arbitrarily small derivative errors near the compact set $\overline W_4\cap\operatorname{supp}\rho$, where $F$ is smooth. On $V$ define
\[
 h=F\quad\hbox{on }W_3,\qquad
 h=\rho F+(1-\rho)g\quad\hbox{on }W_4.
\]
These definitions agree on the overlap. The function $h$ is smooth. Where $\rho$ vanishes it is the section $g$; where $\rho$ can be nonzero, stability applied to the smooth section $F$ makes $h$ a section if the derivative errors of $g-F$ are sufficiently small. Moreover, $|h-F|<\delta$.

Define $f_1=\max\{f,h\}$ on $V$ and $f_1=f$ outside $V$. Near $L_1\cup L_2$ we have $h>f$, and near the boundary of $V$ we have $h<f$. Thus this formula glues continuously, defines a section by maximum closure, and is smooth near $L_1\cup L_2$. Also $\|f_1-f\|_{C^0}<\xi+\delta$. Near $L_1$, it equals $f+\xi$, so the previous smoothness is preserved.

Take finitely many compact sets $L_2$ of this kind covering $M$, and repeat the step, starting with $L_1$ empty and then taking the union of the previously treated compact sets. At every step choose $\xi$ and $\delta$ small enough that the sum of the finitely many uniform errors is below the prescribed tolerance. The final section is smooth everywhere and has the required uniform error.
\end{proof}

\Needspace{5\baselineskip}
\begin{lemma}\label{lem:homogeneous-shear-regularization}
In a holomorphic bundle chart, use an affine projective coordinate $w=v'/v_r$, where $v=(v',v_r)$ and $v_r\ne0$. Suppose that a positive continuous complex homogeneous function $Q$ is strictly plurisubharmonic and locally strongly convex in the fibres on this cone. Put $q(z,w)=Q(z,(w,1))$. For a nonnegative smooth approximate identity $\eta_\delta(s,t)$ of integral one, define
\[
 q_\delta(z,w)=\int q(z-s,w-t)\eta_\delta(s,t)\,\mathrm ds\,\mathrm dt,
 \qquad Q_\delta(z,v)=|v_r|^2q_\delta(z,v'/v_r).
\]
On every relatively compact projective subchart, sufficiently small $\delta$ gives a smooth positive complex homogeneous function with the same strict plurisubharmonicity and local strong fibre convexity. It converges uniformly to $Q$ on compact normalized cone subsets, and in every finite derivative norm wherever $Q$ is already smooth.
\end{lemma}
\begin{proof}
Complex homogeneity gives the identity
\[
 Q_\delta(z,v)
 =\int Q(z-s,T_tv)\eta_\delta(s,t)\,\mathrm ds\,\mathrm dt,
 \qquad T_t(v',v_r)=(v'-v_rt,v_r).
\]
For fixed $(s,t)$ the map $(z,v)\mapsto(z-s,T_tv)$ is holomorphic and invertible, and its fibre restriction is real linear. On a compact normalized cone subset, all these maps for sufficiently small $(s,t)$ stay in a fixed compact subset of the original cone. Their differentials and inverses are uniformly bounded there. A positive local lower bound for the Levi form of $Q$, interpreted through quadratic barriers, therefore gives a positive lower bound for every pullback and for their average. The same argument with real fibre quadratic barriers preserves local strong convexity. Only fibre segments contained in the local cone neighbourhood are used here.

Ordinary convolution makes $q_\delta$ smooth, and the displayed extension makes $Q_\delta$ smooth on the cone and complex homogeneous. Uniform convergence, derivative convergence on compact smooth regions, and positivity for small $\delta$ follow from the corresponding convolution properties and the positive minimum of $Q$ on compact normalized subsets.
\end{proof}

\begin{theorem}\label{thm:homogeneous-convex-regularization}
Let $Q:F\to[0,+\infty)$ be continuous, positive on $F^\times$, complex homogeneous of degree two, strictly plurisubharmonic on $F^\times$, and locally strongly convex in the fibres. There are continuous complex homogeneous functions $Q_j$, smooth on $F^\times$, strictly plurisubharmonic there and locally strongly convex in the fibres, which converge uniformly to $Q$ on the unit sphere bundle of any fixed smooth Hermitian metric. Each $p_j=\sqrt{Q_j}$ is a smooth complex norm.
\end{theorem}
\begin{proof}
Fix a smooth Hermitian squared norm $G$ and let $M=\mathbb P(F)$ denote the bundle of complex lines. The quotient $m=Q/G$ is a positive continuous function on the compact manifold $M$. Define a sheaf $\mathcal S$ on $M$ by requiring that the corresponding function $Gm$ on the cone be positive, strictly plurisubharmonic, and locally strongly convex in the fibres.

This sheaf is closed under maxima. Indeed, multiplication by $G>0$ commutes with maxima. Near a point of the cone, subtract a common positive quadratic barrier from both functions. Their maximum remains plurisubharmonic; on sufficiently small fibre balls it remains convex. Positivity is also preserved.

For stability on a compact subset $K$ of $M$, restrict to normalized vectors projecting to $K$, and to a fixed compact annulus around them. Finitely many bundle charts provide positive common lower bounds for both barriers. If $u$ is a sufficiently small smooth function on $M$ in $C^2$ near $K$, the Levi form and real fibre Hessian of $Gu$ are smaller than half these bounds on that restricted annulus. The perturbed function is still positive and has both strict margins. Homogeneity transports the local conditions to the rest of the cone over $K$. This proves the stated $C^\infty$ stability.

On a sufficiently small projective coordinate chart, \cref{lem:homogeneous-shear-regularization} gives semilocal approximation, including derivative control on compact subsets where $m$ is smooth. Divide its regularizations by $G$ to obtain local smooth sections of $\mathcal S$. By \cref{lem:compact-sheaf-regularization}, there are global smooth positive sections $m_j\to m$ uniformly. Put $Q_j=Gm_j$ off zero and $Q_j=0$ on the zero section. Homogeneity gives continuity at zero, uniformly over the compact base, and gives the asserted convergence on a fixed unit sphere bundle.

For completeness, local fibre convexity away from zero implies convexity on the whole fibre. The restriction to a line avoiding zero is locally convex and hence convex. On a real line through zero, homogeneity gives $Q_j(tv)=t^2Q_j(v)$ for real $t$, which is convex including at zero. Thus the set $B_x=\{v:Q_j(v)\leq1\}$ is convex, balanced and absorbing. Its gauge is $\sqrt{Q_j}$. If $v,w$ are nonzero, convexity applied to $v/p_j(v)$ and $w/p_j(w)$, with weights proportional to $p_j(v)$ and $p_j(w)$, gives the triangle inequality. The cases with a zero vector follow directly. Positivity, smoothness off zero and complex homogeneity give the other properties of a smooth complex norm.
\end{proof}

\begin{proposition}\label{prop:weak-homogeneous-convex-replacement}
Suppose that $F$ admits a smooth positive complex homogeneous function $H$ of degree two that is strictly plurisubharmonic on $F^\times$. If a positive continuous complex homogeneous function $Q$ of degree two is plurisubharmonic on $F^\times$ and locally strongly convex in the fibres, then $F$ admits a smooth complex norm with strictly plurisubharmonic square on $F^\times$.
\end{proposition}
\begin{proof}
For sufficiently small $\tau>0$, set $Q_\tau=Q+\tau H$. It is positive, complex homogeneous and strictly plurisubharmonic. On a compact normalized annulus, local strong fibre convexity of $Q$ gives a positive uniform margin in finitely many frames. The smooth real fibre Hessian of $H$ is uniformly bounded there. Choose $\tau$ so that its possible negative part consumes less than half the margin. Homogeneity then gives local strong fibre convexity everywhere off zero. Apply \cref{thm:homogeneous-convex-regularization} to $Q_\tau$ and take one of its smooth regularizations. Its square root is the required smooth complex norm.
\end{proof}

\begin{corollary}\label{cor:continuous-convex-norm-obstruction}
For every holomorphic structure in the neighbourhood of \cref{thm:convex-norm-deformation-obstruction}, there is no positive continuous complex homogeneous function of degree two that is both strictly plurisubharmonic off zero and locally strongly convex in the fibres. If that structure admits a function $H$ as in \cref{prop:weak-homogeneous-convex-replacement}, there is also no positive continuous complex homogeneous function of degree two that is plurisubharmonic off zero and locally strongly convex in the fibres.
\end{corollary}
\begin{proof}
The first assertion follows from \cref{thm:homogeneous-convex-regularization}, which would produce a smooth complex norm with plurisubharmonic square, contrary to \cref{thm:convex-norm-deformation-obstruction}. For the second assertion use \cref{prop:weak-homogeneous-convex-replacement} and the same obstruction.
\end{proof}

The second assertion requires the specified strictly plurisubharmonic homogeneous function $H$. The argument does not assert that this function exists for every holomorphic bundle or follows from any positivity hypothesis on the base. The regularization does not give a deformation obstruction for continuous fibre norms without a strong convexity margin.
